\section{\benchmark Design}
\label{sec:benchmark-design}

\begin{figure*}[t]
\centering
\begin{tikzpicture}[rb diagram,
  pipeline block/.style={rb block, text width=2.6cm, minimum height=2.1cm},
  pipeline prompt/.style={rb block, text width=4.1cm, minimum height=1cm},
  pipeline oracle/.style={rb edge, dashed}]

\node[rb heading, anchor=west] at (0,1.4) {(a) Dataset construction};
\node[pipeline block] (settings) at (1.5,0)
  {\textbf{Settings}\\Methods, guards, parameters\\Comments, variables, loops};
\node[pipeline block] (java) at (4.9,0)
  {\textbf{Program generation}\\Method tree\\Java class; fixed execution};
\node[pipeline block] (questions) at (8.3,0)
  {\textbf{Oracle and questions}\\Executable calls\\Source--target pairs};
\node[pipeline block, text width=4.3cm, draw=rbInk, line width=0.95pt,
  fill=modelblue!5] (dataset) at (12.9,0)
  {\textbf{\benchmark: final dataset}\\Java classes and reachability questions\\Boolean labels and executable-graph oracles};
\foreach \from/\to in {settings/java,java/questions,questions/dataset}
  \draw[rb edge] (\from.east) -- (\to.west);

\node[rb heading, anchor=west] at (0,-1.85) {(b) Model evaluation};
\node[pipeline prompt] (co) at (4.6,-2.7)
  {\textbf{CO: Class-only}\\Class; entry context and determinism implicit};
\node[pipeline prompt] (ee) at (4.6,-4.25)
  {\textbf{EE: Explicit execution}\\Class + \texttt{Main}; determinism explicit};
\node[pipeline block, text width=2.35cm] (model) at (9.15,-3.475)
  {\textbf{LLM inference}\\For each variant:\\Boolean answer and trace};
\node[pipeline block, text width=3cm] (scoring) at (13.1,-3.475)
  {\textbf{Independent scoring}\\Final-answer correctness\\Trace correctness};

% Only the class and question enter either model prompt.
\draw[rb edge] (dataset.south west) -- (10.575,-1.25)
  -- (-0.2,-1.25) -- (-0.2,-2.7) -- (co.west);
\node[rb annotation] at (6.05,-1.25) {Same class and question};
\draw[rb edge] (-0.2,-2.7) -- (-0.2,-4.25) -- (ee.west);
\draw[rb edge] (co.east) -- (7.55,-2.7) -- (7.55,-3.475) -- (model.west);
\draw[rb edge, -] (ee.east) -- (7.55,-4.25) -- (7.55,-3.475);
\draw[rb edge] (model.east) -- (scoring.west);

% Labels and executable graphs are reserved for evaluation.
\draw[pipeline oracle] (dataset.south) -- (12.9,-1.55)
  -- (13.1,-1.55) -- (scoring.north);
\node[rb annotation, anchor=west, align=left] at (13.25,-1.75)
  {Shared labels\\and oracles};
\end{tikzpicture}
\caption{Construction and evaluation of \benchmark. The construction algorithm produces the final dataset. CO and EE present the same classes and questions with different execution information. Labels and executable-graph oracles are used only for scoring the returned answer and trace.}
\Description{The construction algorithm generates Java classes, records executable calls and selects reachability questions to form ReachabilityBench. The same class and question feed either a Class-only prompt or an Explicit execution prompt with Main and deterministic-execution instructions. Model answers and traces are scored separately against the shared labels and executable-graph oracles, which are not supplied to the model.}
\label{fig:pipeline}
\end{figure*}


\subsection{From a Call Graph to a Reachability Question}

\benchmark starts from a call graph rather than from Java code. The benchmark first generates an ordered graph, then converts each node into a Java method and each edge into a method call. Once the program has been generated, it selects a source method $s$ and a target method $t$. The question asks whether an execution starting from $s$ can reach $t$ through direct or indirect calls.

We represent one benchmark instance with the tuple
\[
    (P, s, t, f, y, \tau, m),
\]
where $P$ is the generated Java program and $f$ identifies the task family. The oracle label $y$ is \true{} when $t$ is executable from $s$, and \false{} otherwise. The oracle trace $\tau$ contains the sequence of methods expected for the question. Finally, $m$ stores the generation parameters and the properties of the graph.

The model receives the complete Java class together with the source and target methods. It must return a JSON object containing two fields:
\begin{itemize}[leftmargin=*]
    \item \texttt{final\_answer}, which contains the Boolean reachability answer; and
    \item \texttt{trace}, which contains the ordered list of methods followed from the source.
\end{itemize}
These two fields are evaluated separately. The final answer shows whether the model reaches the correct decision, while the trace provides information about how it followed the program structure.

Figure~\ref{fig:pipeline} summarizes the complete process. The same generated program is used to build the question, the prompt, and the oracle. After inference, the output is parsed and compared with this oracle. This common origin limits inconsistencies between the program shown to the model and the information used for evaluation.

\subsection{Generating the Call Graph}

The current benchmark uses ordered acyclic trees. This structure allows us to create direct and indirect calls, several branches, and targets that are present in the class but unreachable from the selected source. The order between the children of a node is also important because it defines the order of calls inside the corresponding Java method.

The tree is generated breadth-first under a fixed method budget. Let $u$ be the node currently expanded, and let $L(u)$ and $R(u)$ be its left and right children. When the remaining budget is sufficient, the generator first adds both children:
\[
    u \rightarrow L(u), \qquad u \rightarrow R(u).
\]
It can then add two children below $L(u)$. These additions are sampled independently:
\[
    X_{LL}(u) \sim \mathrm{Bernoulli}(p),
    \qquad
    X_{LR}(u) \sim \mathrm{Bernoulli}(p).
\]
If $X_{LL}(u)=1$, the generator adds $L(u) \rightarrow L(L(u))$. If $X_{LR}(u)=1$, it adds $L(u) \rightarrow R(L(u))$. Figure~\ref{fig:local-tree-rule} shows this local expansion.

\begin{figure}[t]
\centering
\resizebox{\columnwidth}{!}{%
\begin{tikzpicture}[
    >=Latex,
    x=1cm,
    y=1cm,
    method/.style={rb method, minimum width=1.35cm},
    det/.style={rb edge},
    opt/.style={rb optional edge},
    note/.style={font=\scriptsize, align=center}
]
\node[method] (u) at (0,0) {$u$};
\node[method] (l) at (-1.4,-1.0) {$L(u)$};
\node[method] (r) at (1.4,-1.0) {$R(u)$};
\node[method] (ll) at (-2.2,-2.0) {$L(L(u))$};
\node[method] (lr) at (-0.6,-2.0) {$R(L(u))$};
\draw[det] (u) -- node[note, above left] {always} (l);
\draw[det] (u) -- node[note, above right] {always} (r);
\draw[opt] (l) -- node[note, left] {$X_{LL}(u)$} (ll);
\draw[opt] (l) -- node[note, right] {$X_{LR}(u)$} (lr);
\node[rb panel title, text width=4.2cm] at (0,-2.75) {Optional edges are sampled independently with probability $p$.};
\end{tikzpicture}
}
\caption{Local rule used to expand the generated tree. The fixed edges create the main structure, while the optional edges add branches.}
\Description{A node u has deterministic left and right children. The left child has two optional children, each controlled by an independent Bernoulli variable.}
\label{fig:local-tree-rule}
\end{figure}

With this rule, $L(u)$ receives zero, one, or two optional children with probabilities $(1-p)^2$, $2p(1-p)$, and $p^2$. Increasing $p$ increases the number of branches and therefore the number of possible negative questions. The final level can be truncated when the method budget is reached. All labels and traces are computed from this final tree. We use $p=0.65$ to generate enough positive and negative questions; the sensitivity to this value remains to be studied.

\subsection{Generating the Java Program}

After the tree has been generated, each node becomes one Java method. An edge $u \rightarrow v$ inserts a call to $v$ inside the body of $u$. Calls are inserted in the order defined by the tree. In contrast, method declarations are shuffled in the final class. The position of a method in the source file therefore does not reveal the order in which it is executed.

The generator can also add comments, local variables, parameters, loops, and different declaration orders. These elements change the code that is presented to the model, but they do not change reachability unless they are used in a generated condition. In the current benchmark, calls used to define reachability are not placed inside loops.

\subsection{Introducing Runtime Conditions}

The two task families are produced from the same general pipeline, but they differ in the way calls are inserted. In unguarded reachability, every generated call is placed directly in the method body. If execution reaches this statement, the call is executed.

In condition-aware reachability, some calls are placed inside generated \texttt{if} statements. The generator fixes the values used in the conditions and can therefore determine whether each call is executed. The model must reproduce this evaluation from the Java code. It is not sufficient to identify that a call is visible in the method body.

This distinction creates two call graphs. The syntactic graph $G_{\mathit{syn}}$ contains all calls visible in the source code. The executable graph $G_{\mathit{exec}}$ contains only calls whose conditions allow them to execute. The two graphs are equal for unguarded programs. For condition-aware programs, $G_{\mathit{exec}}$ can exclude edges from $G_{\mathit{syn}}$. When a model follows one of these excluded edges, we classify the behavior as \emph{syntactic-call bias}.

\subsection{Building the Oracle Trace}

The oracle is computed after the complete Java program has been generated. It follows executable calls in their source order, starting from the method given in the question. For a positive question, the trace stops when the target is reached for the first time. For a negative question, it continues until the complete region reachable from the source has been explored. Condition-aware traces use $G_{\mathit{exec}}$ and ignore calls whose conditions are false.

This convention means that the canonical trace is not always a shortest path. Execution can visit an earlier sibling branch before reaching the target. Similarly, a negative trace must cover all methods reachable from the source. A model can therefore find a valid path to the target and return the correct Boolean answer without reproducing the complete canonical trace.

\begin{figure}[t]
\centering
\resizebox{\columnwidth}{!}{%
\begin{tikzpicture}[
    >=Latex,
    method/.style={rb method, minimum width=1.05cm},
    edge/.style={rb edge},
    visit/.style={font=\scriptsize, fill=white, inner sep=1pt}
]
\node[method] (a) at (0,0) {A};
\node[method] (b) at (-1.4,-1.0) {B};
\node[method] (c) at (1.4,-1.0) {C};
\node[method] (d) at (-2.1,-2.0) {D};
\node[method] (e) at (-0.7,-2.0) {E};
\node[method] (f) at (0.7,-2.0) {F};
\node[method] (g) at (2.1,-2.0) {G};
\draw[edge] (a) -- node[visit, left] {1} (b);
\draw[edge] (a) -- node[visit, right] {5} (c);
\draw[edge] (b) -- node[visit, left] {2} (d);
\draw[edge] (b) -- node[visit, right] {3} (e);
\draw[edge] (c) -- node[visit, left] {6} (f);
\draw[edge] (c) -- node[visit, right] {7} (g);
\draw[->, line width=1pt, modelblue] ($(a.east)+(0.05,0.05)$) .. controls (3.0,-0.4) and (3.0,-2.0) .. node[visit, right] {positive: A to G} ($(g.east)+(0.05,0)$);
\draw[rb optional edge, line width=1pt] ($(c.west)+(0,0.05)$) .. controls (0.2,-1.0) and (-0.3,-1.5) .. node[visit, above] {negative: C to E} ($(e.north east)+(0,0.03)$);
\end{tikzpicture}
}
\caption{Canonical traversal of a generated call graph. A positive trace includes the calls executed before the target, while a negative trace covers all methods reachable from the source.}
\Description{A tree-shaped call graph with ordered traversal labels and two example reachability queries.}
\label{fig:canonical-traversal}
\end{figure}

We represent the required exploration with a signed canonical trace length. Let $\tau=\langle m_0,\ldots,m_k\rangle$, where $m_0=s$. A positive question has length $+k$, and $m_k=t$. A negative question has length $-k$, where $k$ is the number of steps required to explore the source-reachable region. The sign indicates whether the question is positive or negative, while the magnitude represents the amount of exploration. It is not a shortest-path distance.

\subsection{Configuration and Validation}

The final dataset contains \texttt{[N\_INSTANCES]} questions with context sizes \texttt{[CONTEXT\_SIZES]} and canonical trace lengths \texttt{[DISTANCE\_RANGE]}. Positive and negative labels are balanced at the retained trace lengths. Each released instance contains the Java program, source and target methods, label, canonical trace, task family, class size, graph properties, and generated-condition metadata.

Before an instance is included, the validator compiles its Java class and checks that the generated calls correspond to the intended tree. It then recomputes reachability in $G_{\mathit{exec}}$, verifies the outcome of each generated condition, and checks the canonical trace. As a result, the oracle is validated against the program that is shown to the model rather than only against the state of the generator.
